\section{Limitations} \label{sec:limitations}

The physics in \skye\ models a fully-ionized multi-component quantum ion plasma, quantum and relativistic ideal electrons with non-ideal electron-electron interactions, and ideal radiation.
These components carry with them limitations.
\skye\ is not applicable in the limit of nuclear densities or temperatures: ions are treated as charged point particles and all nuclear interactions are ignored.
Several finite-temperature, composition-dependent, hot nuclear matter EOSs 
have been developed for this regime, including those based on nonrelativistic Skyrme parametrizations
\citep{lattimer_1991_aa,schneider_2017_aa}, variational approaches \citep{togashi_2017_aa} 
and relativistic mean fields \citep{sugahara_1994_aa,shen_1998_aa,typel_2010_aa,fattoyev_2010_aa,steiner_2013_aa}.


Along similar lines at low temperatures and densities, where $T \lesssim 10^{5}\,\mathrm{K}$ and $\rho \lesssim 10^3\,\mathrm{g\,cm^{-3}}$, our ion-ion interaction term becomes large and negative, resulting in unphysical results such as negative entropy.
This reflects the fact that matter is not fully ionized in this limit.
In reality bound states form, reducing the mean ion charge and so reducing the ion-ion interactions.
For very low densities this results in an ideal gas with a different mean molecular weight.
Several EOSs have been developed for this regime, including those based on 
free energy minimization \citep{Saumon1995,Irwin2004},
cluster activity expansions \citep{rogers_1974_aa,rogers_1981_aa,Rogers2002},
cluster viral expansions \citep{omarbakiyeva_2015_aa,ballenegger_2018_aa},
density-functional theory molecular dynamics \citep{militzer_2013_aa,2014ApJS..215...21B}, 
path integral Monte Carlo \citep{militzer_2001_aa}, 
quantum Monte Carlo \citep{mazzola_2018_aa},
Feynman-Kac path integral representations \citep{alastuey_2020_aa}, and
asymptotic expansions \citep{alastuey_2012_aa}.
Using these EOSs in stellar evolution calculations typically requires 
pre-tabulating results for fixed compositions due to the computational cost of solving for ionization equilibrium.

In principle partial ionization could be included in \skye\ in a variety of ways.
For instance we could add terms accounting for electron-ion interactions, but unfortunately we are not aware of robust prescriptions for the interaction free energy $F_{j-e}$ in this limit.
The challenge is that existing prescriptions are based on perturbation expansions~\citep{1961ApJ...134..669S,Potekhin2010}, but these break down well before the formation of bound states~\citep{1976PhRvA..14..816G}.
Variational approaches seem more promising in this limit, but are more computationally expensive to implement because they involve minimizing the free energy with respect to a variational parameter~\citep{1976PhRvA..14..816G}.
The same is true for direct solutions to the Saha equation, which are generally quite expensive.

A further limitation concerns our understanding of high density
quantum melts.  The physics is not as well understood as for lower
densities or higher temperatures. We think this is a fruitful area 
for further study, particularly given that the quantum melt line
\skye\ currently predicts disagrees with calculations based
on the Lindemann criterion~\citep{1993ApJ...414..695C,PhysRevB.18.3126,PhysRevLett.76.4572}.

Putting these limitations together, we recommend that \skye\ not be used for densities above $0.1 \rho_{j,\rm nuclear} \approx A_j 10^{13}\mathrm{g\,cm^{-3}}$, where $A_j$ is the number of baryons per ion, or for temperatures above the proton rest mass-energy $T_{\rm QCD} \approx 10^{13}\,\mathrm{K}$.
We further recommend that \skye\ not be used in the joint limit $T < T_{j}^{\rm ion}$ and $\rho < \rho_{j}^{\rm ion}$.
Here $T_{j}^{\rm ion}$ is the temperature above which a dilute gas is fully ionized.
Neglecting degeneracy factors, we may solve for this using the Saha equation
\begin{align}
	\frac{n_{j,Z_j}}{n_{j,Z_j-1}} = \frac{2 n_{Q,e}}{n_e} e^{-\psi_{f,j} / k_{\rm B} T},
\end{align}
where $\psi_{f,j}$ is the final ionization potential of a species of charge $Z_{j}$, and  $n_{j,Z}$ is the number density of fully ionized ions of species $j$ and charge $Z$.
As a rough heuristic we require $n_{j,Z_j} > 10 n_{n_{j,Z_j-1}}$ to ensure that full ionization is a good approximation.
With this we find
\begin{align}
	k_{\rm B} T_{j}^{\rm ion} \approx \frac{\psi_{f,j}}{\frac{3}{2} \ln (T_{j}^{\rm ion}/10^4\,\mathrm{K}) - \ln (Z_j \rho / A_j \mathrm{g\,cm^3}) - 7}.
\end{align} 
If we approximate $\psi_{f,j} \approx \mathrm{Ry} Z_j^2$ we then find
\begin{align}
	T_{j}^{\rm ion} \approx \frac{10^5\,\mathrm{K} Z_j^2}{\frac{3}{2} \ln (T_{j}^{\rm ion}/10^4\,\mathrm{K}) - \ln (Z_j \rho / A_j \mathrm{g\,cm^3}) - 7}.
\end{align} 
For densities below that of pressure ionization this typically gives $T_{j}^{\rm ion} \approx 10^4\,\mathrm{K} Z_j^2$.
Along similar lines, $\rho_{j}^{\rm ion}$ is the density above which a low-temperature system is fully ionized, given approximately by~\citep{doi:10.1098/rspa.1938.0073}
\begin{align}
	\rho_{j}^{\rm ion} &= \frac{3 m_{j}}{\pi \sqrt{2}} \left(\frac{\psi_{f,j}}{e^2 a_0}\right)^{3/2}\\
	&\approx 3 \frac{m_j}{m_p} Z_j^{3} \mathrm{g\,cm^{-3}}\approx 3 A_j Z_j^{3} \mathrm{g\,cm^{-3}},
\end{align}
where $a_0 = \hbar^2 / m_e e^2$ is the Bohr radius.
For mixtures of ions we recommend averaging $\rho_j^{\rm ion}$ and $T_j^{\rm ion}$ weighted by number density to determine the appropriate limits.
Finally, we recommend caution in interpreting results in the quantum melt limit, which occurs in the joint limit of $\rho > (A_j/12)^4 (Z_j/6)^6 10^9 \mathrm{g\,cm^{-3}}$ and $T < (A_j/12)(Z_j/6)^4 10^7\mathrm{K}$.
